% ============================================================
% SECTION 10 — DATABASE ARCHITECTURE
% ============================================================

\section{Database Architecture}

\begin{tcolorbox}[summarybox={\iconDatabase[1.5] The Indisputable Source of Truth}]
\color{textdark}
In RecoverAI V2, the relational database (PostgreSQL) is the \textbf{exclusive source of truth}. While Redis is used for ephemeral task queuing, any data not durably committed to PostgreSQL does not exist in the financial reality of the system.

The schema is defined in \filepath{app/models/db\_models.py} using SQLAlchemy 2.0. It is designed to be fully normalized and highly resilient to concurrent mutations.
\end{tcolorbox}

\vspace{0.8cm}

% --- ENTITY RELATIONSHIP DIAGRAM ---
\subsection{Entity Relationship Diagram}

The following diagram illustrates the core models and their foreign key relationships.

\vspace{0.5cm}

\begin{center}
\begin{tikzpicture}[
    node distance=1cm and 2cm,
    entity/.style={rectangle, draw=primary, thick, fill=bglight, minimum width=3.5cm, minimum height=1cm, rounded corners=4pt, font=\small\bfseries, drop shadow=black!5, align=center},
    attribute/.style={font=\scriptsize, align=left, text=textdark},
    arrow/.style={-{Stealth[scale=1.1]}, thick, draw=textdark},
    rel/.style={font=\scriptsize\itshape, text=textmuted}
]

    % Core Entities
    \node[entity] (txn) {Transaction};
    \node[attribute, below=0cm of txn] {
        \code{id} (PK) \\
        \code{amount} (Integer) \\
        \code{status} \\
        \code{recovery\_status} \\
        \code{refund\_status}
    };
    
    \node[entity, right=3cm of txn] (attempt) {RecoveryAttempt};
    \node[attribute, below=0cm of attempt] {
        \code{id} (PK) \\
        \code{transaction\_id} (FK) \\
        \code{status} \\
        \code{version} (Integer, OCC)
    };
    
    \node[entity, below=3.5cm of txn] (audit) {AuditLog};
    \node[attribute, below=0cm of audit] {
        \code{id} (PK) \\
        \code{transaction\_id} (FK) \\
        \code{action} \\
        \code{details} (JSON)
    };
    
    \node[entity, below=3.5cm of attempt] (webhook) {WebhookEvent};
    \node[attribute, below=0cm of webhook] {
        \code{id} (PK) \\
        \code{status} \\
        \code{retry\_count} \\
        \code{last\_attempt\_at}
    };
    
    \node[entity, right=2.5cm of attempt] (idemp) {IdempotencyRecord};
    \node[attribute, below=0cm of idemp] {
        \code{key} (PK) \\
        \code{request\_hash} \\
        \code{status}
    };

    % Relationships
    \draw[arrow] (txn.east) -- node[above, rel] {1 : N} (attempt.west);
    \draw[arrow] (txn.south) -- node[left, rel] {1 : N} (audit.north);

\end{tikzpicture}
\end{center}

\vspace{0.8cm}

% --- CORE ENTITIES ---
\subsection{Core Entity Definitions}

\subsubsection{1. \code{Transaction}}
The root of the data model. Represents the original payment attempt made by the customer.
\begin{itemize}
    \item \textbf{Financial Math:} Amounts are strictly stored as \textbf{Minor Units} (integers). For example, \$50.00 is stored as \code{5000} to prevent floating-point rounding errors during financial execution.
    \item \textbf{Dual Statuses:} \code{recovery\_status} tracks the aggregate state of the recovery process (\code{NOT\_STARTED}, \code{EXECUTING}, \code{SUCCEEDED}), while \code{refund\_status} independently tracks refund actions.
\end{itemize}

\vspace{0.4cm}

\subsubsection{2. \code{RecoveryAttempt}}
A child record representing a single, specific effort to recover a transaction.
\begin{itemize}
    \item \textbf{1:N Relationship:} A single \code{Transaction} can have multiple \code{RecoveryAttempts} if the policy engine permits multiple retries.
    \item \textbf{Optimistic Lock:} Contains the \code{version} integer field. Every time the attempt transitions state, this version is atomically incremented.
\end{itemize}

\vspace{0.4cm}

\subsubsection{3. \code{AuditLog}}
An append-only log providing a cryptographic-like history of what the system did and why.
\begin{itemize}
    \item \textbf{Compliance:} Records every LLM diagnosis, policy decision, and gateway execution.
    \item \textbf{Immutability:} By design, the application never updates or deletes an \code{AuditLog} row once inserted.
\end{itemize}

\vspace{0.4cm}

\subsubsection{4. \code{IdempotencyRecord}}
The shield for the API layer. 
\begin{itemize}
    \item \textbf{Scoping:} Keys are generated dynamically based on the merchant's hashed API key combined with the client-provided header (\code{pay\_[api\_hash]\_[idemp\_key]}).
    \item \textbf{Hash Validation:} Stores a \code{request\_hash} of the payload. If a client reuses an idempotency key but alters the payload (e.g., changes the amount from \$50 to \$100), the system rejects it as a collision.
\end{itemize}

\vspace{0.4cm}

\subsubsection{5. \code{WebhookEvent}}
A durable inbox for gateway callbacks.
\begin{itemize}
    \item \textbf{Ingestion Separation:} Webhook payloads are inserted into this table in a raw, unparsed format before being picked up by Celery. This allows the HTTP listener to return \code{200 OK} instantly.
    \item \textbf{Batch 5.3.1 Retries:} Contains \code{retry\_count} and \code{last\_attempt\_at} to prevent infinite loops when parsing fails permanently.
\end{itemize}

\vspace{0.8cm}

% --- SQLITE VS POSTGRESQL ---
\subsection{The SQLite vs. PostgreSQL Divide}

\begin{tcolorbox}[warningbox={The Danger of SQLite in Financial Systems}]
RecoverAI V2 was built and tested primarily using SQLite. While SQLite is fantastic for local development, it possesses a structural limitation that is catastrophic for financial queues: \textbf{it does not support row-level locking}.
\end{tcolorbox}

When the application executes \code{db.query(Model).with\_for\_update()}, SQLAlchemy passes the \code{SELECT ... FOR UPDATE} command to the database engine.
\begin{itemize}
    \item \textbf{PostgreSQL:} Locks the specific row. If Worker 2 tries to read it, it is forced to wait until Worker 1 commits.
    \item \textbf{SQLite:} Silently ignores the \code{FOR UPDATE} clause. Worker 2 reads the row simultaneously.
\end{itemize}

\textbf{Mitigation applied in V2:} Because SQLite ignored pessimistic locks, the engineering team was forced to implement \textbf{Optimistic Concurrency Control} (\code{UPDATE ... WHERE status = 'X' AND version = 'Y'}) for all state transitions (Batch 4.6). OCC relies on basic atomic writes, which both SQLite and PostgreSQL support, ensuring the system is safe regardless of the underlying engine.

\begin{tcolorbox}[successbox={Production Readiness}]
Despite the OCC mitigations making the code safe on SQLite, the system is explicitly designed to be deployed against \textbf{PostgreSQL} in production to leverage true pessimistic locks in the reconciliation sweepers (preventing CPU waste when multiple sweepers conflict).
\end{tcolorbox}
